% ============================================================
% Script de soutenance — Myriam (diapos 1 à 4)
% Compilation : pdflatex script_myriam.tex
% ============================================================
\documentclass[11pt,a4paper]{article}

\usepackage[utf8]{inputenc}
\usepackage[T1]{fontenc}
\usepackage[french]{babel}
\usepackage[margin=2.4cm]{geometry}
\usepackage{amsmath,amssymb}
\usepackage{booktabs}
\usepackage{parskip}

\newcommand{\diapo}[2]{\section*{Diapo #1 --- #2}}
\newenvironment{comprendre}%
  {\par\medskip\noindent\hrulefill\par\nopagebreak\textbf{Pour bien comprendre :}\itshape\par}%
  {\par\nopagebreak\noindent\hrulefill\par\medskip}

\title{\textbf{Script de soutenance --- Myriam}\\
\large Effet tunnel et temps de traversée d'une barrière de potentiel\\
\normalsize Diapos 1 à 4 --- environ 5 minutes}
\author{}
\date{Soutenance juin 2026 --- oraux du 22 au 26 juin}

\begin{document}
\maketitle

\subsection*{Répartition des rôles}

\begin{center}
\begin{tabular}{llc}
\toprule
Qui & Diapos & Durée visée \\
\midrule
\textbf{Myriam} & 1--4 : problématique, paquet d'ondes, états stationnaires & $\approx$ 5 min \\
Sheryne & 5--8 : méthode numérique, validation, premiers résultats & $\approx$ 4,5 min \\
Mattéo & 9--12 : effet Hartman, études paramétriques, conclusion & $\approx$ 5,5 min \\
\bottomrule
\end{tabular}
\end{center}

Total : 15 minutes, puis 10 minutes de questions. Le texte ci-dessous est à dire
à l'oral ; les phrases \textbf{en gras} sont les messages essentiels, à ne jamais
sauter même si le temps presse.

% ------------------------------------------------------------
\diapo{1}{Page de titre (30 secondes)}

Bonjour à toutes et à tous. Nous allons vous présenter notre projet numérique de
physique moderne, qui porte sur \textbf{l'effet tunnel et le temps de traversée
d'une barrière de potentiel} par un paquet d'ondes gaussien à une dimension.

Je suis Myriam Bensaid, et je travaille avec Sheryne Ouarghi et Mattéo Knopes.
Je vais d'abord poser le problème et les outils théoriques ; Sheryne présentera
ensuite la méthode numérique et les premiers résultats ; et Mattéo terminera par
l'étude détaillée du temps de traversée et la conclusion.

% ------------------------------------------------------------
\diapo{2}{Problématique (1 minute)}

En mécanique classique, la situation est simple : une particule d'énergie $E$
qui arrive sur une barrière de hauteur $V_0$ supérieure à $E$ rebondit, toujours.
En mécanique quantique, c'est différent : la particule a une probabilité non
nulle de passer de l'autre côté. C'est l'effet tunnel.

Notre question centrale va un cran plus loin : \textbf{combien de temps met une
particule quantique pour traverser une barrière qu'elle n'a classiquement pas le
droit de franchir ?} Cette question paraît innocente, mais elle est encore
débattue aujourd'hui, et nous verrons qu'elle réserve des surprises.

Notre démarche suit cinq étapes, que vous voyez à l'écran : construire un paquet
d'ondes gaussien ; déterminer les états stationnaires de la barrière ; résoudre
numériquement l'équation de Schrödinger dépendante du temps ; mesurer le temps de
traversée et le comparer à la théorie ; et enfin étudier l'influence des
paramètres de la barrière.

Tout le projet est en unités adimensionnées, $\hbar = m = 1$ : cela allège les
équations sans rien changer à la physique.

\begin{comprendre}
Pourquoi la question du temps est-elle difficile ? Parce que sous la barrière,
l'onde n'oscille pas : elle s'amortit (on dit qu'elle est « évanescente »). Il
n'y a donc pas de vitesse de propagation évidente à l'intérieur, et il faut
\emph{définir} ce qu'on appelle « temps de traversée ». Nous choisirons une
définition opérationnelle : le temps qui sépare le passage du pic du paquet à
l'entrée et à la sortie de la barrière.
\end{comprendre}

% ------------------------------------------------------------
\diapo{3}{Le paquet d'ondes gaussien (1,5 minute)}

Première étape : représenter une particule. Une onde plane unique ne convient
pas, car elle est complètement délocalisée, présente partout dans l'espace. Pour
obtenir une particule localisée, on superpose une infinité d'ondes planes : c'est
l'intégrale à l'écran, où chaque onde de vecteur d'onde $k$ est pondérée par un
poids gaussien $g(k)$ centré sur une valeur moyenne $k_0$.

\textbf{Ce paquet a quatre propriétés clés.} Un : son centre avance à la vitesse
de groupe $v_g = k_0$ dans nos unités --- c'est elle qui joue le rôle de la
vitesse de la particule. Deux : les oscillations internes avancent à la vitesse
de phase, qui vaut la moitié de $v_g$. Trois : sa largeur spectrale vaut
$\Delta k = 1/(2\sqrt{a_{\rm wp}})$, où $a_{\rm wp}$ contrôle la largeur
spatiale ; \textbf{plus le paquet est étroit en position, plus il est large en
vecteur d'onde} --- c'est exactement le principe d'indétermination de Heisenberg.
Quatre : le paquet s'étale au cours du temps, car la relation de dispersion
$\omega = k^2/2$ n'est pas linéaire : les composantes rapides prennent de
l'avance sur les lentes.

Sur la figure de droite, vous voyez les parties réelle et imaginaire du paquet à
l'instant initial : des oscillations sous une enveloppe gaussienne.

\begin{comprendre}
Retenez surtout $\Delta k$ : toute la fin de l'exposé repose dessus. Le paquet
n'a pas \emph{une} énergie mais \emph{une plage} d'énergies de largeur liée à
$\Delta k$. Or la barrière ne traite pas toutes ces composantes de la même
façon : c'est ce « filtrage » qui expliquera les écarts entre nos mesures et la
théorie onde plane.
\end{comprendre}

% ------------------------------------------------------------
\diapo{4}{États stationnaires de la barrière (2 minutes --- diapo indispensable)}

Deuxième outil, et c'est un point que le jury attend : les états stationnaires.
On cherche les solutions d'énergie $E$ fixée, avec $E < V_0$, pour une barrière
rectangulaire de hauteur $V_0$ et de largeur $a$.

L'espace se découpe en trois zones. Avant la barrière : une onde incidente et une
onde réfléchie d'amplitude $R$. Après la barrière : une onde transmise
d'amplitude $T$. Et \textbf{dans la barrière, le vecteur d'onde devient
imaginaire} : la solution n'est plus oscillante mais exponentielle, en
$e^{\pm\kappa x}$, avec $\kappa = \sqrt{2m(V_0 - E)}/\hbar$. C'est l'onde
évanescente : l'onde ne se propage pas sous la barrière, elle s'y amortit.

Pour relier les trois zones, on impose la continuité de $\psi$ et de sa dérivée
en $x = 0$ et en $x = a$ : quatre équations, quatre inconnues. On en tire
l'amplitude de transmission $T(k)$ affichée dans l'encadré.

Deux conséquences. D'abord, \textbf{$|T|^2$ n'est jamais nul : la transmission
est toujours possible}, c'est la signature mathématique de l'effet tunnel.
Ensuite, en régime opaque, c'est-à-dire $\kappa a \gg 1$, la transmission décroît
comme $e^{-2\kappa a}$ : \textbf{une sensibilité exponentielle à la largeur et à
la hauteur de la barrière}. C'est ce qui rend l'effet tunnel si sélectif --- et
si utile, par exemple dans le microscope à effet tunnel.

Enfin, notez que $T$ est un nombre complexe : il a un module, mais aussi une
phase $\varphi_T$. \textbf{Cette phase contient l'information sur le temps de
traversée} : Mattéo s'en servira pour construire la prédiction théorique.

Je laisse maintenant la parole à Sheryne pour la résolution numérique.

\begin{comprendre}
Le raccordement est une question classique de jury : il faut savoir le refaire
au tableau. La logique : dans chaque zone l'équation de Schrödinger est une
équation différentielle linéaire d'ordre 2 à coefficients constants, donc la
solution générale est une combinaison de deux exponentielles ; les conditions de
continuité aux deux interfaces fixent les amplitudes relatives. La continuité de
$\psi'$ vient du fait que le potentiel reste fini (il n'y a un saut de dérivée
que pour un potentiel en pic de Dirac).
\end{comprendre}

% ------------------------------------------------------------
\clearpage
\section*{Vue d'ensemble de l'exposé (pour pouvoir répondre sur tout)}

\textbf{Le fil rouge en quatre phrases.} Nous voulons mesurer le temps que met un
paquet d'ondes à traverser une barrière par effet tunnel. Pour cela, nous avons
besoin de deux outils : le paquet gaussien (partie de Myriam) et la transmission
$T(k)$ des états stationnaires (Myriam aussi). Sheryne montre que notre solveur
Crank--Nicolson est fiable (norme conservée, temps libre $\tau_0$ retrouvé à
4\,\% près) et que le paquet se scinde en une partie réfléchie dominante et une
petite partie transmise. Mattéo montre alors le résultat principal :
\textbf{le paquet tunnel sort plus tôt que le paquet libre ($\tau_t < \tau_0$),
c'est l'effet Hartman}, et son comportement face à la barrière (saturation
avec $a$, diminution avec $V_0$). Il évoque en perspective l'influence des
caractéristiques du paquet lui-même ($k_0$, $a_{\rm wp}$), explorée mais
écartée du diaporama par souci de temps.

\textbf{Les nombres à connaître par c\oe ur} (cas de référence $k_0 = 5$,
$E = 12{,}5$, $V_0 = 15$, $a = 1$, $a_{\rm wp} = 1$) :

\begin{center}
\begin{tabular}{lcc}
\toprule
Grandeur & Numérique & Théorique \\
\midrule
Temps libre $\tau_0$ & 0{,}192 & 0{,}200 \\
Temps tunnel $\tau_t$ & 0{,}112 & 0{,}168 (Hartman) \\
Transmission & $P_{\rm trans} = 8{,}1\times10^{-2}$ & $|T(k_0)|^2 = 2{,}5\times10^{-2}$ \\
Saturation de $\tau_t^{\rm th}$ avec $a$ & --- & $\approx 0{,}179$ \\
\bottomrule
\end{tabular}
\end{center}

\section*{Questions probables du jury --- éléments de réponse}

\textbf{Refaire le raccordement des états stationnaires au tableau.}
Voir l'encadré de la diapo 4 ci-dessus : solutions par zone, continuité de
$\psi$ et $\psi'$ en $0$ et $a$, résolution du système $4\times4$.

\textbf{Pourquoi Crank--Nicolson plutôt qu'Euler explicite ?}
Parce qu'il est \textbf{unitaire} (la norme, donc la probabilité totale, est
conservée exactement par le schéma) et \textbf{inconditionnellement stable},
alors qu'Euler explicite fait exploser la norme. C'est la partie de Sheryne,
mais chacun doit savoir le dire.

\textbf{Avez-vous étudié l'influence de $k_0$ et $a_{\rm wp}$ ? Les temps
négatifs violent-ils la causalité ?}
Oui, nous l'avons fait, mais écarté du diaporama final par souci de temps
(perspective de la conclusion). Non, pas de violation : le pic du paquet
transmis n'est pas un objet matériel qui se propage, il est \emph{reconstruit}
à partir des composantes rapides du spectre, que la barrière favorise
fortement. Aucune information ne va plus vite que prévu ; c'est un effet de
filtrage spectral.

\textbf{Pourquoi la probabilité transmise mesurée dépasse-t-elle $|T(k_0)|^2$ ?}
Parce que le paquet n'est pas monochromatique : il contient des composantes
$k > k_0$ bien mieux transmises, et $|T|^2$ croît très vite avec $k$. Dans la
limite quasi monochromatique ($a_{\rm wp}$ grand, $\Delta k$ petit), la mesure
convergerait vers $|T(k_0)|^2$ --- une piste creusée mais pas gardée dans le
diaporama.

\textbf{Effet Hartman et superluminalité apparente : que se passe-t-il pour
$a \to \infty$ ?}
$\tau_t^{\rm th}$ sature alors que la distance croît : la « vitesse » apparente
$a/\tau_t$ deviendrait arbitrairement grande. Mais en pratique $|T|^2$ s'effondre
en $e^{-2\kappa a}$ et la mesure du pic perd son sens bien avant : il n'y a pas
de transport d'information superluminal.

\end{document}
